\documentclass{article}
\usepackage{amsmath}
\usepackage{french-logic}
\begin{document}
\begin{theorem}
  Test \cax and \CMr and \cmr and $\deland p$ and $p \defby q$ and \Kfour and \Rup.
\end{theorem}
\begin{proof}
  A proof with \mprule and \close and \infr and \qed, plus $\mlnot p \mland q$.
  Ladder and adjacency: $\M,\omega\trues\ought A\to\may B\land C$ and $\logic\proves C$.
\end{proof}
\begin{definition}
  Uses $\tuple{a,b}$ and $\af p$ and \EN and $\Rlog$ and $\pzero$ and $\cnecs$ and $\langd$
  and $\langle w\rangle$.
\end{definition}
\begin{hilbertlist}
  \item x \by{RR}
\end{hilbertlist}
Arg test $\truthset{ZZ}$ and $\proofsetl{YY}$ and \hyp{QQ} and \parent{PP}.
\begin{tabular}{c}\toprule x\\ \midrule y\\ \bottomrule\end{tabular}
\begin{proofsketch}
  Sketchy.
\end{proofsketch}
% citation commands: \tcite/\pcite must colour exactly as \cite,
% locator or not (after/syntax/tex.lua routes them through texCmdRef).
\cite{ckA} \cite[p.~5]{ckB} \tcite{ckC} \tcite[p.~7]{ckD}
\pcite{ckE} \pcite[\S2]{ckF} \textcite{ckG} \poscite{ckH}
\end{document}
